\subsection{惯性与参考系}
\begin{tabularx}{\columnwidth}{XX}
  牛顿第一定律 & 不受外力时物体保持静止或匀速直线运动 \\
  惯性量度 & 质量是惯性的量度，质量越大惯性越强 \\
  惯性系定义 & 相对惯性系静止或匀速直线运动的参考系 \\
  非惯性系特点 & 参考系有加速度，需引入惯性力辅助分析 \\
\end{tabularx}

\begin{formulanotebox}
\begin{itemize}[leftmargin=1.5em]
  \item 牛顿第一定律揭示“力是改变运动状态的原因”，不是维持运动的原因。
  \item 非惯性系中若直接用牛顿第二定律，通常需补加惯性力：$\vec F_\text{惯}=-m\vec a_\text{系}$。
  \item “惯性大”不代表“速度大”，而代表“运动状态更难被改变”。
\end{itemize}
\end{formulanotebox}

\begin{keypointbox}
伽利略斜面实验通过“理想化外推”得到无阻力极限：若无外力，速度保持不变，这一思想是建立经典力学体系的关键起点。
\end{keypointbox}

\begin{casetypebox}[title={\strut\textcolor{white}{\bfseries 常见题型：惯性定律理解辨析}}]
% TODO: 力与运动关系、惯性与速度混淆辨析题型待补充
\end{casetypebox}

\begin{casetypebox}[title={\strut\textcolor{white}{\bfseries 常见题型：参考系选择}}]
% TODO: 惯性系与非惯性系判定、参考系变换下运动描述题型待补充
\end{casetypebox}

\begin{examplebox}
% TODO: 惯性与参考系专题例题待补充
\end{examplebox}

\begin{examplesolutionbox}
% TODO: 对应例题解析待补充
\end{examplesolutionbox}
